\begin{table}[t]
\caption{Five families of methods for estimating prevalence from a classifier}
\label{tab:families}
\centering
\footnotesize
\begin{tabular}{@{}p{0.31\columnwidth}p{0.43\columnwidth}c@{}}
\toprule
Family (labels in the tables) & How it estimates & Each group \\
\midrule
\raggedright 1. Classify and count (CC) & Label each sentence with one of the four rules of Section~\ref{sec:rules} and count the labels. & No \\
\addlinespace
\raggedright 2. Probability averages (PCC) & Average the probabilities, raw or after Platt scaling on validation data. & No \\
\addlinespace
\raggedright 3. Adjusted counts (ACC, PACC, EM, ACC-G) & Undo the classifier's error rates, measured on validation data, in each group for ACC-G. & ACC-G \\
\addlinespace
\raggedright 4. Slope correction (SC) & Divide each contrast by the shrinkage measured on validation contrasts of the same kind. & No \\
\addlinespace
\raggedright 5. Per-group estimates (annotated share, PPI, Platt per group) & Use each group's annotated documents, alone or to correct the classifier's estimate. & Yes \\
\bottomrule
\multicolumn{3}{@{}p{0.96\columnwidth}@{}}{\scriptsize Each group: whether the method needs annotated documents from every group compared.}
\end{tabular}
\end{table}
